%
% 2026-07-04-screen-state-machine.tex
%
% Z Specification for the XBoing screen / mode state machine.
%
% Models the code AS IS: the flat screen mode, the attract-cycle
% transition constraint, the shared inner sub-screen systems, and the
% shared resources (game-active flag, required-block grid, high-score
% board selection) that the mode-enter and input handlers read/write.
%
% Peer-reviewed against the current code on 2026-07-04; the board-load,
% board-selection, and game-active-lifecycle facts below are grounded
% in cited file:line evidence.
%

\documentclass[a4paper,10pt,fleqn]{article}
\usepackage[margin=1in]{geometry}
\usepackage{fuzz}
\usepackage[colorlinks=true,linkcolor=blue,citecolor=blue,urlcolor=blue]{hyperref}

% Long \texttt identifiers (e.g. game_callbacks.c:750--759) do not
% hyphenate; let TeX stretch interword space to avoid overfull boxes.
\setlength{\emergencystretch}{4em}

\begin{document}

\title{XBoing Screen State Machine: A Z Specification}
\author{Formal Model of the Flat Mode Machine and its Shared Resources}
\date{July 2026}
\hypersetup{
  pdftitle={XBoing Screen State Machine: A Z Specification},
  pdfauthor={xboing-c},
  pdfsubject={Z Specification},
  pdfcreator={fuzz/probcli}
}
\maketitle

\tableofcontents
\newpage

\section{Introduction}

The XBoing screen machine is \emph{flat}.  \texttt{struct sdl2\_state}
(\texttt{src/sdl2\_state.c}) holds a single \textsf{current} mode
(\texttt{SDL2ST\_*}), plus \textsf{previous} and a \textsf{saved} slot
for dialogue push/pop.  There is no enclosing ``attract'' super-state:
\textsf{Intro}, \textsf{Instruct}, \textsf{Demo}, \textsf{Preview},
\textsf{Keys}, \textsf{KeysEdit} and \textsf{Highscore} are all
peer top-level modes.  The ``attract loop'' is a designated
\emph{subset} with a constrained transition order (\texttt{attract\_cycle[]}),
i.e.\ a constraint on a flat machine, not a nesting.

A second, orthogonal nesting axis does exist but is out of scope here:
while \textsf{current} sits on an attract mode, that mode's driving
\emph{system} runs its own animation phase (\texttt{intro\_state\_t},
\texttt{demo\_state\_t}, \dots) held inside \texttt{ctx->intro} /
\texttt{ctx->demo} / \texttt{ctx->keys}, not inside \texttt{sdl2\_state}.
Three systems are each shared by two modes (\texttt{intro\_system}
  drives \textsf{Intro}/\textsf{Instruct}; \texttt{demo\_system} drives
  \textsf{Demo}/\textsf{Preview}; \texttt{keys\_system} drives
\textsf{Keys}/\textsf{KeysEdit}).

Three resources are \emph{shared} between the attract screens and real
gameplay: the \texttt{game\_active} flag, the required-block grid
(\texttt{block\_system\_still\_active}), and the high-score tables
(\texttt{hs\_global} / \texttt{hs\_personal}, chosen by
\texttt{highscore\_request\_type}).  The observed defects all arise
because a transition reads a shared resource that a prior transition
left in the wrong state.  This model names the safety invariants the
handlers must preserve and shows which operations break them.

\section{Free Types}

The outer modes.  Legacy \textsf{BallWait}/\textsf{Wait} (defined but
never assigned) are omitted.

\begin{zed}
  Mode ::= mNone | mPresents | mIntro | mInstruct \\
  \quad~ | mDemo | mPreview | mKeys | mKeysedit \\
  \quad~ | mHighscore | mBonus | mGame | mPause \\
  \quad~ | mDialogue | mEdit
\end{zed}

The shared inner systems (descriptive: the nesting axis noted above).

\begin{zed}
  System ::= sysNone | sysIntro | sysDemo | sysKeys
\end{zed}

Which high-score table the display is bound to
(\texttt{highscore\_request\_type}).

\begin{zed}
  Board ::= bGlobal | bPersonal
\end{zed}

fuzz has no built-in boolean; \textsf{ZBOOL} avoids the ProB/B keyword
clash with \textsf{BOOL}.

\begin{zed}
  ZBOOL ::= ztrue | zfalse
\end{zed}

\section{Screen Sets and the Attract Cycle}

The attract screens form the auto-cycle; \textsf{Presents} is not part
of it.

\begin{axdef}
  attractModes : \power Mode
  \where
  attractModes = \{ mIntro, mInstruct, mDemo, mPreview, \\
  \quad~ mKeys, mKeysedit, mHighscore \}
\end{axdef}

Attract screens on which \textsf{Space} starts a game directly.
\textsf{Intro} (goes to \textsf{Instruct}) and \textsf{Highscore}
(conditional) are handled separately (\texttt{game\_input.c:399--423}).

\begin{axdef}
  spaceStartModes : \power Mode
  \where
  spaceStartModes = \{ mInstruct, mDemo, mPreview, mKeys, mKeysedit \}
\end{axdef}

The cycle order used by the C key and each screen's finish-timer
auto-advance (\texttt{attract\_cycle[]}, \texttt{game\_callbacks.c:54--68}).

\begin{axdef}
  attractNext : Mode \pfun Mode
  \where
  attractNext = \{ mIntro \mapsto mInstruct, \\
    \quad~ mInstruct \mapsto mDemo, \\
    \quad~ mDemo \mapsto mKeys, \\
    \quad~ mKeys \mapsto mKeysedit, \\
    \quad~ mKeysedit \mapsto mPreview, \\
    \quad~ mPreview \mapsto mHighscore, \\
  \quad~ mHighscore \mapsto mIntro \}
\end{axdef}

The shared system driving each attract screen (two-to-one sharing).

\begin{axdef}
  driver : Mode \pfun System
  \where
  driver = \{ mIntro \mapsto sysIntro, \\
    \quad~ mInstruct \mapsto sysIntro, \\
    \quad~ mDemo \mapsto sysDemo, \\
    \quad~ mPreview \mapsto sysDemo, \\
    \quad~ mKeys \mapsto sysKeys, \\
  \quad~ mKeysedit \mapsto sysKeys \}
\end{axdef}

\section{State}

The flattened screen state.  \textsf{blocksLoaded} abstracts
\texttt{block\_system\_still\_active} (required blocks present).
\textsf{globalHasData} / \textsf{personalHasData} record whether each
table held rows after the startup read
(\texttt{game\_init.c:327--335}); on an unprivileged install (no
\texttt{/var/games}) \textsf{globalHasData} is \textsf{zfalse} while
\textsf{personalHasData} may be \textsf{ztrue}.  \textsf{boardChoice}
is \texttt{highscore\_request\_type}.

\begin{schema}{Screen}
  mode : Mode \\
  gameActive : ZBOOL \\
  blocksLoaded : ZBOOL \\
  boardChoice : Board \\
  globalHasData : ZBOOL \\
  personalHasData : ZBOOL
\end{schema}

The machine boots into \textsf{Presents}: \texttt{ctx} is
\texttt{calloc}'d so \textsf{gameActive} is \textsf{zfalse}
(\texttt{game\_init.c:231}); no level is loaded; both boards are read
from disk (\texttt{game\_init.c:327--335}) so the data flags reflect
what was on disk --- left unconstrained here because they are
platform-dependent; and \textsf{boardChoice} defaults to
\textsf{bGlobal} (\texttt{game\_init.c:340}).

\begin{schema}{Init}
  Screen'
  \where
  mode' = mPresents \\
  gameActive' = zfalse \\
  blocksLoaded' = zfalse \\
  boardChoice' = bGlobal
\end{schema}

\section{Safety Invariants (the intended contract)}

Stated as observer schemas, deliberately \emph{not} folded into
\textsf{Screen}: in Z the state predicate constrains every
$\Delta\textsf{Screen}$, so folding an invariant in would silently
\emph{disable} the violating transitions rather than expose them (probcli
then reports a spurious deadlock).  Instead each contract is machine-checked
by model-checking toward its \emph{negation} as a probcli goal:
\begin{quote}
\texttt{probcli spec.tex -model\_check -goal "<negated invariant>"}
\end{quote}
A reachable state matching the negation is a counter-example proving the
defect; ``\textsf{No counter example found. ALL states visited}'' after a
fix proves the guard closes it.

A game may only run once its level's required blocks are on the grid;
otherwise the per-tick rule check reads an empty grid and declares the
level complete.

\begin{schema}{SafeGame}
  Screen
  \where
  mode = mGame \implies blocksLoaded = ztrue
\end{schema}

The high-score screen must not display an \emph{empty} board when a
populated one is available: if it is bound to an empty global board,
the personal board must also be empty (nothing better to show).

\begin{schema}{SafeHighscore}
  Screen
  \where
  (mode = mHighscore \\
    \quad~ \land boardChoice = bGlobal \\
  \quad~ \land globalHasData = zfalse) \\
  \implies personalHasData = zfalse
\end{schema}

\textsf{game\_active} may legitimately stay true on the game-over
\textsf{Highscore} display (it gates score submission), but it must
not propagate to the \emph{other} attract screens; otherwise the
\textsf{Highscore} \textsf{Space} routing misfires.

\begin{schema}{SafeAttract}
  Screen
  \where
  (mode \in attractModes \land mode \neq mHighscore) \\
  \implies gameActive = zfalse
\end{schema}

\section{Operations (as coded)}

\subsection{Boot: Presents auto-advances to Intro}

The machine boots into \textsf{Presents}; when the credits animation
finishes, \texttt{presents\_cb\_on\_finished}
(\texttt{game\_callbacks.c:596--600}) transitions unconditionally to
\textsf{Intro}, entering the attract cycle.  \textsf{Presents} is not
part of \texttt{attract\_cycle[]}.

\begin{schema}{Boot}
  \Delta Screen
  \where
  mode = mPresents \\
  mode' = mIntro \\
  gameActive' = gameActive \\
  blocksLoaded' = blocksLoaded \\
  boardChoice' = boardChoice \\
  globalHasData' = globalHasData \\
  personalHasData' = personalHasData
\end{schema}

\subsection{Attract auto-advance and the C key}

Both the finish-timer (\texttt{highscore\_cb\_on\_finished} and peers)
and the C key advance the cycle.  Crucially, this transition
\emph{preserves} \textsf{gameActive} --- the auto-advance from
\textsf{Highscore} does not clear it (\texttt{game\_callbacks.c:750--759}).
Entering \textsf{Highscore} binds the display to \textsf{boardChoice}
but performs no disk read and no global-to-personal fallback on the
attract path (\texttt{game\_modes.c:1139--1142}).

\begin{schema}{AttractAdvance}
  \Delta Screen
  \where
  mode \in \dom attractNext \\
  mode' = attractNext(mode) \\
  (mode' = mIntro \implies gameActive' = zfalse) \\
  (mode' \neq mIntro \implies gameActive' = gameActive) \\
  blocksLoaded' = blocksLoaded \\
  (mode' = mHighscore \land gameActive = zfalse \\
    \quad~ \land globalHasData = zfalse \land personalHasData = ztrue \\
  \quad~ \implies boardChoice' = bPersonal) \\
  (\lnot (mode' = mHighscore \land gameActive = zfalse \\
    \quad~ \land globalHasData = zfalse \land personalHasData = ztrue) \\
  \quad~ \implies boardChoice' = boardChoice) \\
  globalHasData' = globalHasData \\
  personalHasData' = personalHasData
\end{schema}

\subsection{Starting a game --- level load succeeds}

\texttt{start\_new\_game} (\texttt{game\_modes.c:90--166}) sets
\textsf{gameActive} and, on a successful load, populates the grid via
the \texttt{on\_add\_block} callback.  This is the correct path.

\begin{schema}{StartNewGameOk}
  \Delta Screen
  \where
  mode \in spaceStartModes \\
  mode' = mGame \\
  gameActive' = ztrue \\
  blocksLoaded' = ztrue \\
  boardChoice' = boardChoice \\
  globalHasData' = globalHasData \\
  personalHasData' = personalHasData
\end{schema}

\subsection{Starting a game --- level load fails silently}

If the level file cannot be resolved or parsed
(\texttt{game\_modes.c:128--136}) \texttt{start\_new\_game} only warns
to stderr and still enters \textsf{Game} with the grid left cleared.

\begin{schema}{StartNewGameFail}
  \Delta Screen
  \where
  mode \in spaceStartModes \\
  mode' = mGame \\
  gameActive' = ztrue \\
  blocksLoaded' = zfalse \\
  boardChoice' = boardChoice \\
  globalHasData' = globalHasData \\
  personalHasData' = personalHasData
\end{schema}

\subsection{Space on the title}

\begin{schema}{IntroSpace}
  \Delta Screen
  \where
  mode = mIntro \\
  mode' = mInstruct \\
  gameActive' = gameActive \\
  blocksLoaded' = blocksLoaded \\
  boardChoice' = boardChoice \\
  globalHasData' = globalHasData \\
  personalHasData' = personalHasData
\end{schema}

\subsection{Space on the high-score screen}

The \textsf{Space} handler routes on \textsf{gameActive} and clears it
here (\texttt{game\_input.c:411--423}): a live game-over returns to
\textsf{Intro}; otherwise it starts a game.

\begin{schema}{HighscoreSpaceReturn}
  \Delta Screen
  \where
  mode = mHighscore \\
  gameActive = ztrue \\
  mode' = mIntro \\
  gameActive' = zfalse \\
  blocksLoaded' = blocksLoaded \\
  boardChoice' = boardChoice \\
  globalHasData' = globalHasData \\
  personalHasData' = personalHasData
\end{schema}

\begin{schema}{HighscoreSpaceStart}
  \Delta Screen
  \where
  mode = mHighscore \\
  gameActive = zfalse \\
  mode' = mGame \\
  gameActive' = ztrue \\
  blocksLoaded' = ztrue \\
  boardChoice' = boardChoice \\
  globalHasData' = globalHasData \\
  personalHasData' = personalHasData
\end{schema}

\subsection{Game over}

\texttt{game\_rules\_ball\_died} (\texttt{game\_rules.c:294--296})
leaves \textsf{gameActive} true and shows \textsf{Highscore}.

\begin{schema}{GameOver}
  \Delta Screen
  \where
  mode = mGame \\
  mode' = mHighscore \\
  gameActive' = ztrue \\
  blocksLoaded' = blocksLoaded \\
  (globalHasData = zfalse \land personalHasData = ztrue \\
  \quad~ \implies boardChoice' = bPersonal) \\
  (\lnot (globalHasData = zfalse \land personalHasData = ztrue) \\
  \quad~ \implies boardChoice' = bGlobal) \\
  globalHasData' = globalHasData \\
  personalHasData' = personalHasData
\end{schema}

\subsection{The per-tick rule check}

Every \textsf{Game} tick, \texttt{game\_rules\_check}
(\texttt{game\_rules.c:330--340}) sends the machine to \textsf{Bonus}
the instant no required blocks remain.  It is \emph{unconditional} ---
no guard that a level was actually loaded.

\begin{schema}{GameRulesToBonus}
  \Delta Screen
  \where
  mode = mGame \\
  blocksLoaded = zfalse \\
  mode' = mBonus \\
  gameActive' = gameActive \\
  blocksLoaded' = blocksLoaded \\
  boardChoice' = boardChoice \\
  globalHasData' = globalHasData \\
  personalHasData' = personalHasData
\end{schema}

\subsection{Bonus completes: next level loads}

On bonus completion \texttt{bonus\_cb\_on\_finished} advances to the next
level (\texttt{game\_rules\_next\_level}) and re-enters \textsf{Game} with
the new level's blocks loaded (\texttt{game\_callbacks.c:671--678}).
Modelled so \textsf{Bonus} is not a terminal (deadlocking) mode.

\begin{schema}{BonusToGame}
  \Delta Screen
  \where
  mode = mBonus \\
  mode' = mGame \\
  blocksLoaded' = ztrue \\
  gameActive' = gameActive \\
  boardChoice' = boardChoice \\
  globalHasData' = globalHasData \\
  personalHasData' = personalHasData
\end{schema}

\section{Conjectures: where the invariants break}

Each defect is a reachable state that violates one safety schema.  A
model checker (probcli) explores \textsf{Init} plus the operations and
reports the trace; the prose states the trace and the missing guard.

\subsection*{Defect 2 --- empty high-score board in the attract cycle
  (PROVEN, FIXED)}

The boards \emph{are} read at startup, so ``never loaded'' is not the
cause.  On an unprivileged install \textsf{globalHasData} =
\textsf{zfalse} and \textsf{personalHasData} = \textsf{ztrue}, and
\textsf{boardChoice} = \textsf{bGlobal} from \textsf{Init}.  Goal-directed
model checking of the negated \textsf{SafeHighscore} reached
\textsf{mode} = \textsf{mHighscore}, \textsf{gameActive} = \textsf{zfalse}
still bound to the empty global board by pure attract-cycling
($Init \semi Boot \semi AttractAdvance^{*}$), the machine-checked
counter-example for the reported symptom.

\textbf{Fix.} \texttt{mode\_highscore\_enter}
(\texttt{src/game\_modes.c}) now falls back to \textsf{bPersonal} on the
attract path (\texttt{!game\_active}) when the global board is empty
(\texttt{highscore\_io\_count(\&hs\_global) == 0}) and the personal board
has scores.  The game-over path keeps its own submit-based fallback
(\texttt{game\_modes.c:1134}).  Both entry paths are modelled here in
\textsf{AttractAdvance} and \textsf{GameOver}; re-running the goal now
reports \emph{No counter example found. ALL states visited} (138 states),
so the fixed operations cannot reach the violating state.  Regression:
\texttt{test\_highscore\_attract\_falls\_back\_to\_personal}
(\texttt{tests/test\_integration\_modes.c}) plus
\texttt{highscore\_io\_count} unit tests.

\subsection*{Defect 3 --- first Space advances the loop (PROVEN, FIXED)}

In a \emph{pure} attract cycle \textsf{gameActive} is always
\textsf{zfalse} (set only by \textsf{StartNewGame*}, cleared only by
the \textsf{Highscore} \textsf{Space} handler), so
\textsf{IntroSpace} (\textsf{Intro} $\to$ \textsf{Instruct}) is the
\emph{intended} PR~\#172 behavior, not a bug.  The genuine violation
requires a prior game over: goal-directed model checking of the negated
\textsf{SafeAttract} found the 7-step trace
$\dots \semi GameOver \semi AttractAdvance$ carrying
\textsf{gameActive} = \textsf{ztrue} out of \textsf{Highscore} into
\textsf{Intro} (\textsf{SafeAttract} violated at \textsf{Intro}); the
next time the cycle reaches \textsf{Highscore},
\textsf{HighscoreSpaceReturn} fires and the first \textsf{Space} goes to
\textsf{Intro} instead of starting a game.

\textbf{Fix.} Every attract advance out of \textsf{Highscore} --- the
finish-timer (\texttt{highscore\_cb\_on\_finished}), the C cycle key
(\texttt{game\_input.c}), and the \textsf{Space} return --- lands on
\textsf{Intro}, since \texttt{attract\_next(HIGHSCORE) == INTRO}.  A
single clear in \texttt{mode\_intro\_enter} (\texttt{src/game\_modes.c})
therefore closes all three leak paths at one site; modelled here as
\textsf{AttractAdvance} setting \textsf{gameActive'} = \textsf{zfalse}
whenever \textsf{mode'} = \textsf{mIntro}.  Re-running the goal now
reports \emph{No counter example found. ALL states visited} (90 states),
and no non-\textsf{Highscore} attract mode is reachable with
\textsf{gameActive} = \textsf{ztrue}.  Regression:
\texttt{test\_highscore\_autoadvance\_clears\_game\_active} (timer path)
and \texttt{test\_highscore\_c\_cycle\_clears\_game\_active} (C key).
See ADR-055.

\subsection*{Defect 1 --- Space lands in Bonus with an empty grid}

\textsf{StartNewGameOk} establishes \textsf{blocksLoaded}, so the model
exonerates the \emph{happy-path} \textsf{Space}-to-game edge.  The
violation is \textsf{StartNewGameFail}: a silent level-load failure
(\texttt{game\_modes.c:128--136}) enters \textsf{Game} with
\textsf{blocksLoaded} = \textsf{zfalse}, and the very next
\textsf{GameRulesToBonus} tick drops to \textsf{Bonus}
(\textsf{SafeGame} violated).  \textbf{Missing guard:}
\texttt{start\_new\_game} must not enter \textsf{Game} on a failed
load, and/or \texttt{game\_rules\_check} must require a
positive ``level populated'' precondition rather than trusting the
grid.

\section{Test-coverage gap}

The mode suite (\texttt{tests/test\_integration\_modes.c})
force-transitions to each mode and ticks for \emph{no-crash} only; it
asserts nothing about the shared resources.  None of the three
invariants is guarded:

\begin{itemize}
  \item No test drives $Init \semi AttractAdvance^{*}$ to
    \textsf{Highscore} with an empty global board and asserts the display
    falls back to a populated personal board (\textsf{SafeHighscore}).
  \item No test drives \textsf{GameOver} then auto-cycles and asserts
    \textsf{gameActive} is cleared before the next attract screen
    (\textsf{SafeAttract}), nor that the first \textsf{Space} after an
    auto-cycled game-over starts a game.
  \item No test asserts entering \textsf{Game} on a failed level load
    does not immediately fall to \textsf{Bonus} (\textsf{SafeGame}); the
    one bonus-fixture test deliberately uses an \emph{empty} grid,
    exercising the opposite.
\end{itemize}

Each invariant should become a checked postcondition of the
corresponding handler and a regression test.

\end{document}
